\scanpage{220}
\margin{De Rham P.D.}
\textcircled{3} Sur $\Gamma^*M\otimes\Lambda^*M$ \deleted{(algèbre bigraduée)}, il y a un seul op. $d$ qui soit \deleted{a) $A$ linéaire} a) une \inserted{anti-}dérivation \unsure{de|à} \inserted{$A$-algèbre graduée} \deleted{$\Gamma^*(M)$ de degrés pairs et $\Lambda M$ de degré impair} \inserted{anti-dérivation de degré $1$} b) compatible avec les puiss. divisées sur $(\Gamma^*M)^+$ i.e. $d(x^{(i)})=x^{(i-1)}dx$, et c) $d$ est nulle sur $\Lambda^1M$, et sur $\Gamma^1M$ c'est l'isom can $\Gamma^1M\longrightarrow\Lambda^1M$. Alors $d$ satisfait $d^2=0$, et applique $\Gamma^p\otimes\Lambda^q$ dans $\Gamma^{p-1}\otimes\Lambda^{q+1}$. Il est $\Lambda^*M$-linéaire. C'est aussi du ``complexe de De Rham à puissances divisées'' de la $A$-algèbre à p.d. $\Gamma^*M$ (avec son idéal $\Gamma^+M$).
\margin{$?\ \left\{\vphantom{\begin{matrix}1\\2\\3\end{matrix}}\right.$}
Si $M$ est plat sur $A$, alors on a (graduant par le degré de $\Lambda^*$)
\[
H^*_{DR\,\mathrm{p.d.}}(\Gamma^*(M)\otimes\Lambda^*M)\simeq
\begin{cases}0&\text{si }i\neq0\\
\underset{\substack{\Vert\\\Gamma^0\otimes\Lambda^0}}{A}&\text{si }i=0
\end{cases}
\]

\margin{Koszul P.D.}
\textcircled{4} Sur $\Gamma^*M\otimes\Lambda^*M$, il y a un seul op. $\delta$ qui soit

a) dérivation de \inserted{$A$-}algèbre graduée b) compatible aux p.d. sur $\Gamma^+(M)$ $(d(x^{(i)})=x^{(i-1)}dx\text{ si }x\in\Gamma^+)$.

c) nulle sur $\Gamma^1M$ \deleted{(donc sur $\Gamma^*M$ \unsure{entier})} et induisant sur $\Lambda^1M$ l'isom $\Lambda^1M\longrightarrow\Gamma^1M$. \deleted{Si $M$ \deleted{libre} \inserted{plat} sur} \deleted{Alors et si $A$ de car. $0$,} \inserted{alors $\delta^2=0$, il applique} $\Gamma^p\otimes\Lambda^q$ dans $\Gamma^{p+1}\otimes\Lambda^{q-1}$. \inserted{est $\Gamma^*M$-linéaire.}
Écrivant les degrés \deleted{$\unsure{M}$} en bas, si $M$ est plat sur $A$ de car. $0$, on a
\margin{(inutile)}
\[
H_i^{\mathrm{Kos\ p.d.}}(\Gamma^*M\otimes\Lambda^*M)\simeq
\begin{cases}0&\text{si }i\neq0\\
\underset{\substack{\Vert\\\Gamma^0\otimes\Lambda^0}}{A}&\text{si }i=0
\end{cases}
\]
Si $A$ de car $p>0$, $M$ plat, on doit avoir une histoire d'opérations de Cartier ----
\margin{219}
